%% gallery.tex -- one page, one seed; used by tools/gallery.py to render the seed gallery.
%%   lualatex -jobname=build/seed-7 "\def\florilegiumseed{7}%% gallery.tex -- one page, one seed; used by tools/gallery.py to render the seed gallery.
%%   lualatex -jobname=build/seed-7 "\def\florilegiumseed{7}%% gallery.tex -- one page, one seed; used by tools/gallery.py to render the seed gallery.
%%   lualatex -jobname=build/seed-7 "\def\florilegiumseed{7}%% gallery.tex -- one page, one seed; used by tools/gallery.py to render the seed gallery.
%%   lualatex -jobname=build/seed-7 "\def\florilegiumseed{7}\input{gallery.tex}"
\documentclass[11pt,a4paper]{article}
\renewcommand\rmdefault{cmr}\renewcommand\sfdefault{cmss}\renewcommand\ttdefault{cmtt}
\usepackage[T1]{fontenc}
\usepackage{amsmath,amssymb}
\usepackage[a4paper,left=68mm,right=22mm,top=32mm,bottom=30mm,headsep=8mm]{geometry}
\usepackage{fancyhdr}
\usepackage[grown]{florilegium}
\pagestyle{fancy}\fancyhf{}
\fancyhead[R]{\small\textsc{3\quad Procedural ornament}}
\renewcommand\headrulewidth{0pt}
\begin{document}
\setcounter{section}{3}\setcounter{subsection}{1}
\subsection{One seed, one garden}

This page was grown from seed \texttt{\ifx\florilegiumseed\empty page\else\florilegiumseed\fi}.
Nothing in the margin was placed by hand: three stems rise from a root in the corner and
sway up the column, leaves hang at their nodes, and a greedy loop then finds the largest
remaining hole, chooses a motif that fits it and grows a branch to it from the nearest
stem, until only small gaps are left. The same seed always gives the same garden; a
different seed gives a different one. Because the text on this page ends early, the
garden also grows into the space below these lines.
\florilegiumtextend
\end{document}
"
\documentclass[11pt,a4paper]{article}
\renewcommand\rmdefault{cmr}\renewcommand\sfdefault{cmss}\renewcommand\ttdefault{cmtt}
\usepackage[T1]{fontenc}
\usepackage{amsmath,amssymb}
\usepackage[a4paper,left=68mm,right=22mm,top=32mm,bottom=30mm,headsep=8mm]{geometry}
\usepackage{fancyhdr}
\usepackage[grown]{florilegium}
\pagestyle{fancy}\fancyhf{}
\fancyhead[R]{\small\textsc{3\quad Procedural ornament}}
\renewcommand\headrulewidth{0pt}
\begin{document}
\setcounter{section}{3}\setcounter{subsection}{1}
\subsection{One seed, one garden}

This page was grown from seed \texttt{\ifx\florilegiumseed\empty page\else\florilegiumseed\fi}.
Nothing in the margin was placed by hand: three stems rise from a root in the corner and
sway up the column, leaves hang at their nodes, and a greedy loop then finds the largest
remaining hole, chooses a motif that fits it and grows a branch to it from the nearest
stem, until only small gaps are left. The same seed always gives the same garden; a
different seed gives a different one. Because the text on this page ends early, the
garden also grows into the space below these lines.
\florilegiumtextend
\end{document}
"
\documentclass[11pt,a4paper]{article}
\renewcommand\rmdefault{cmr}\renewcommand\sfdefault{cmss}\renewcommand\ttdefault{cmtt}
\usepackage[T1]{fontenc}
\usepackage{amsmath,amssymb}
\usepackage[a4paper,left=68mm,right=22mm,top=32mm,bottom=30mm,headsep=8mm]{geometry}
\usepackage{fancyhdr}
\usepackage[grown]{florilegium}
\pagestyle{fancy}\fancyhf{}
\fancyhead[R]{\small\textsc{3\quad Procedural ornament}}
\renewcommand\headrulewidth{0pt}
\begin{document}
\setcounter{section}{3}\setcounter{subsection}{1}
\subsection{One seed, one garden}

This page was grown from seed \texttt{\ifx\florilegiumseed\empty page\else\florilegiumseed\fi}.
Nothing in the margin was placed by hand: three stems rise from a root in the corner and
sway up the column, leaves hang at their nodes, and a greedy loop then finds the largest
remaining hole, chooses a motif that fits it and grows a branch to it from the nearest
stem, until only small gaps are left. The same seed always gives the same garden; a
different seed gives a different one. Because the text on this page ends early, the
garden also grows into the space below these lines.
\florilegiumtextend
\end{document}
"
\documentclass[11pt,a4paper]{article}
\renewcommand\rmdefault{cmr}\renewcommand\sfdefault{cmss}\renewcommand\ttdefault{cmtt}
\usepackage[T1]{fontenc}
\usepackage{amsmath,amssymb}
\usepackage[a4paper,left=68mm,right=22mm,top=32mm,bottom=30mm,headsep=8mm]{geometry}
\usepackage{fancyhdr}
\usepackage[grown]{florilegium}
\pagestyle{fancy}\fancyhf{}
\fancyhead[R]{\small\textsc{3\quad Procedural ornament}}
\renewcommand\headrulewidth{0pt}
\begin{document}
\setcounter{section}{3}\setcounter{subsection}{1}
\subsection{One seed, one garden}

This page was grown from seed \texttt{\ifx\florilegiumseed\empty page\else\florilegiumseed\fi}.
Nothing in the margin was placed by hand: three stems rise from a root in the corner and
sway up the column, leaves hang at their nodes, and a greedy loop then finds the largest
remaining hole, chooses a motif that fits it and grows a branch to it from the nearest
stem, until only small gaps are left. The same seed always gives the same garden; a
different seed gives a different one. Because the text on this page ends early, the
garden also grows into the space below these lines.
\florilegiumtextend
\end{document}
